% spintent-07.tex
% Guía de ejercicios de geometría: ángulos, áreas y perímetros.
%
% Ejemplo real: guía para estudiantes de octavo básico.
%
% Las lecturas en los comentarios `% doc:` son las de NVDA 2026.2 con
% MathCAT, idénticas en mathml-SE y en mathml-AF: una por cada fórmula,
% en el orden en que aparece. «a» y «b» (en el vector) y «centímetros
% cuadrado» son lectura propia de MathCAT, no del paquete.
%
% Compilar con: lualatex-dev spintent-07.tex
\DocumentMetadata
  {
    lang=es-CL,
    pdfversion=2.0,
    pdfstandard=ua-2,
    tagging=on,
    tagging-setup={math/setup={mathml-SE}},
  }
\documentclass{article}
\usepackage[spanish]{babel}
\usepackage{spanish-mathml-intent, hyperref}
\hypersetup
  {
    pdfdisplaydoctitle = true,
    pdftitle           = {Guía de geometría},
  }

\begin{document}

\section*{Guía \spshort{n.o} 5: ángulos, áreas y perímetros}

Preparada por la \spshort{prof.a} Pérez para el curso de \spshort{8.o} básico.
Use regla y compás cuando se le indique.

\subsection*{Ángulos}

\begin{enumerate}
  % doc: «A B C»
  % doc: «medida del ángulo en A es igual a 40 grados»
  % doc: «medida del ángulo en B es igual a 65 grados»
  % doc: «medida del ángulo en C»
  \item En el triángulo $\spPoly[concept=false]{ABC}$ se cumple que $\spmangle{A} = \spang{40}$ y $\spmangle{B} = \spang{65}$. Calcule $\spmangle{C}$.
  % doc: «A O B»
  % doc: «37 grados 30 minutos»
  \item Mida con un transportador el ángulo $\spangle[concept=false]{AOB}$ y escriba su medida con minutos: $\spang{37:30}$.
  % doc: «90 grados»
  % doc: «medida del ángulo alfa es igual a 35 grados»
  % doc: «medida del ángulo beta»
  \item Dos ángulos son complementarios si suman $\spang{90}$. Si $\spmangle{\alpha} = \spang{35}$, calcule $\spmangle{\beta}$.
  % doc: «A O B»
  \item Dibuje un ángulo recto $\spangle[recto=true,concept=false]{AOB}$ y trace su bisectriz.
\end{enumerate}

\subsection*{Puntos, rectas y segmentos}

\begin{enumerate}
  % doc: «A 1 coma 2»
  % doc: «B 4 coma 2»
  % doc: «C 4 coma 6»
  \item Ubique en el plano cartesiano los puntos $\spPoint[concept=false]{A}(1,2)$, $\spPoint[concept=false]{B}(4,2)$ y $\spPoint[concept=false]{C}(4,6)$.
  % doc: «A B»
  % doc: «B C»
  \item Trace la recta $\sprecta[concept=false]{AB}$ y el rayo $\sprayo[concept=false]{BC}$.
  % doc: «A C»
  % doc: «medida del lado A B es igual a 3 centímetros»
  % doc: «medida del lado B C es igual a 4 centímetros»
  \item Calcule la medida del segmento $\spside[concept=false]{AC}$ sabiendo que $\spmside{AB} = \spunit{3 cm}$ y $\spmside{BC} = \spunit{4 cm}$.
  % doc: «a»
  % doc: «b»
  % doc: «vector A B es igual a 3 coma 0»
  \item Escriba el vector de $A$ a $B$: $\spvec{AB}(3,0)$.
\end{enumerate}

\subsection*{Áreas y perímetros}

\begin{enumerate}
  % doc: «A B C»
  % doc: «6 centímetros»
  % doc: «4 centímetros»
  % doc: «área del triángulo A B C es igual a 12 centímetros cuadrado»
  \item El triángulo $\spPoly[concept=false]{ABC}$ tiene base $\spunit{6 cm}$ y altura $\spunit{4 cm}$. Calcule su área: $\spAfig(ABC) = \spunit{12 cm^{2}}$.
  % doc: «A B C D»
  % doc: «5 centímetros»
  % doc: «perímetro del cuadrado A B C D es igual a 20 centímetros»
  \item Un cuadrado $\spPoly[concept=false]{ABCD}$ tiene lado $\spunit{5 cm}$. Calcule su perímetro: $\spPfig(ABCD) = \spunit{20 cm}$.
  % doc: «8 centímetros»
  % doc: «3 centímetros»
  \item Calcule el área de un rectángulo de lados $\spunit{8 cm}$ y $\spunit{3 cm}$.
  % doc: «n»
  % doc: «perímetro del polígono»
  \item Escriba la fórmula del perímetro de un polígono de $n$ lados iguales: $\spPfig(P)$.
\end{enumerate}

\subsection*{Circunferencia y círculo}

\begin{enumerate}
  % doc: «circunferencia de centro en O y radio r»
  % doc: «radio es igual a 3 centímetros»
  \item Dibuje $\spcirc(O, r)$ con $\sprad = \spunit{3 cm}$.
  % doc: «diámetro es igual a 6 centímetros»
  \item Calcule su diámetro: $\spdiam = \spunit{6 cm}$.
  % doc: «área del círculo con centro en O»
  % doc: «perímetro de la circunferencia con centro en O»
  \item Escriba la fórmula para calcular $\spAcirc[read-sty=circulo](O)$ y $\spPcirc(O)$.
  % doc: «A B»
  % doc: «medida del arco A B»
  \item Marque el arco $\sparc[concept=false]{AB}$ y escriba $\spmarc{AB}$.
\end{enumerate}

\subsection*{Volumen}

\begin{enumerate}
  % doc: «5 centímetros»
  % doc: «volumen es igual a 125 centímetros al cubo»
  \item Un cubo tiene arista $\spunit{5 cm}$. Calcule su volumen: $\spVol = \spunit{125 cm^{3}}$.
\end{enumerate}

% doc: «dos de mayo de dos mil veintisiete»
Entregue esta guía el $\spdate{02/05/2027}$.

\end{document}